\documentclass[10pt,a4paper]{amsart}
\usepackage[margin=19mm]{geometry}
\usepackage{amsmath,amssymb,mathtools}
\usepackage[scr=rsfs,cal=euler]{mathalfa}
\usepackage{xfrac}
\usepackage[unicode,hidelinks,bookmarks=false]{hyperref}
\newcommand{\ie}{i.e.\ }
\newcommand{\cf}{cf.\ }
\newcommand{\resp}{respectively\ }
\newcommand{\Subsection}{\subsection}
\providecommand{\nobreakdash}{\nobreak\hbox{-}}
\setlength{\emergencystretch}{2em}
\multlinegap0pt
\usepackage{amssymb}
\usepackage{xcolor}
\newtheorem{theorem}{Theorem}
\DeclareMathOperator{\I}{Id}
\DeclareMathOperator{\supp}{supp}
\title{Multipliers, $W$-algebras and the growth of generalized polynomial identities}
\author{Fabrizio Martino, Carla Rizzo}
\date{Source: arXiv:2502.12830v3 (CC BY 4.0)}
\begin{document}
\maketitle

\noindent\emph{Source and licence.} Multipliers, $W$-algebras and the growth of generalized polynomial identities. Version arXiv:2502.12830v3 (CC BY 4.0), distributed by arXiv under the Creative Commons Attribution 4.0 International licence, http://creativecommons.org/licenses/by/4.0/, as recorded in the arXiv licence metadata for this exact version and shown on the abstract page https://arxiv.org/abs/2502.12830v3. Full licence notice: \texttt{licenses/arxiv-2502.12830-CC-BY-4.0.txt}.\par\smallskip

\noindent\textit{Editorial scope.} Counted unit: the article's theorem identitadiEk (source0.tex lines 967--974). The article's complete original proof of it is delivered with it (source0.tex lines 975--1021, one \texttt{proof} environment of 47 lines). No mathematics is written, repaired or re-derived: the statement and the proof are the article's own lines, in the article's own notation, and no retained line is substituted or re-flowed.  No further source-local context is needed: every cross-reference of every retained line resolves inside the two delivered spans.  Nothing was omitted from any retained span, so the package carries no omission record.  No retained line contains a citation, so the wrapper carries no bibliography and no bibliography entry.  No retained line contains a figure environment, an image-inclusion command or drawing code, so no image is delivered and the package declares no asset dependency.  Editorial formatting only, in addition: an AMS \texttt{amsart} preamble that restates the article's own declarations and macros used by the retained lines, namely the environments \texttt{theorem} on separate counters, together with the standard packages those lines need.  The retained environments print in this document as Theorem 1 in order of appearance, whereas the article prints the same items with the numbering of its own full document.  The complete native source of the article as distributed in the arXiv e-print for this exact version is bundled unchanged as \texttt{sources/173477/source0.tex}.
\begin{theorem}\label{identitadiEk}
    Let $E^{(k)},$ with $k\geq 1$, be the unital Grassmann algebra $E$ over a field $F$ of characteristic zero regarded as an $E$-algebra, where $E$ acts as the subalgebra $E_k$ by left and right multiplication.
    Then $\I^E(E^{(k)})$ is generated, as $T_E$-ideal, by the polynomials 
    \begin{equation*}
    [x_1,x_2,x_3], \quad [e_i,x_1,x_2], \quad e_jx, \quad xe_j
\end{equation*}
for all $1\leq i\leq k$ and for all $j\geq k+1.$ Moreover, $c^E_n(E^{(k)}) = 2^{n-1}(2^{k+1}-1)$ for all $n\geq 1.$
\end{theorem}

\begin{proof}
    Let $I$ be the $T_E$-ideal generated by the above generalized polynomials. A straightforward computation shows that $I\subseteq \I^E(E^{(k)}).$ Thus in order to prove the opposite inclusion, let $f\in\I^E(E^{(k)})$ be a multilinear generalized polynomial of degree $n$ and suppose, as we may, that in $f$ does not appear any $e_j$ for all $j\geq k+1.$ 

    By the Poincaré-Birkhoff-Witt Theorem and by $[x_1,x_2,e_i]\equiv 0$, $f$ can be written as a linear combination of polynomials of the type
    $$
    x_{c_1}\cdots x_{c_t}v_1v_2\cdots v_me_{l_1}\cdots e_{l_s},
    $$
    where $c_1<\cdots < c_t,$ $l_1<\cdots <l_s,$ $0\leq s\leq k$ and $v_1,\ldots, v_m$ are left-normed commutators containing variables and $e_i$'s. Moreover, since $E$ has unity, without loss of generality in what follows we can disregard the tail $x_{c_1}\cdots x_{c_t}$ and suppose that the variables $x_1,\ldots,x_n$ in $f$ appear exclusively within the commutators.
    
    Now, using Lemma \ref{identitadiEk} and the $W$-trivial polynomials mentioned above, we can write $f$ as linear combination of the polynomials
    \begin{equation}\label{polinomigradopari}
    [x_{1},x_{2}][x_3,x_4]\cdots [x_{2m-1},x_{2m}]e_{l_1}\cdots e_{l_s},
    \end{equation}
    if $n=2m$ is even or 
    \begin{equation}\label{polinomigradodispari}
    [e_{l_1},x_{1}][x_2,x_3]\cdots [x_{2m},x_{2m+1}]e_{l_2}\cdots e_{l_s},
    \end{equation}
    if $n=2m+1$ is odd. In both cases $l_1<\cdots <l_s$ and $0\leq s\leq k.$ This implies that such polynomials generate $P^W_n$ modulo the generalized identities of $E.$

    We claim that they are also linearly independent modulo $\I^E(E^{(k)}).$ To this end, first suppose that $\deg f= 2m$ and write
    $$
    f=  \sum_{L}\alpha_L[x_{1},x_{2}][x_3,x_4]\cdots [x_{2m-1},x_{2m}]e_{l_1}\cdots e_{l_s} \quad(\text{mod }\I^E(E^{(k)})),
    $$
    where $L=\{l_1,\ldots, l_s\}.$ Let $s$ be the smallest integer for which by contradiction there exists $L$ such that $\alpha_L\neq 0.$ Then, by making the evaluation $x_1\mapsto g_1= e_{j_1}\cdots e_{j_t}g'$ and $x_i\mapsto g_i$ for all $2\leq i\leq 2m,$ where $\{j_1,\ldots, j_t\}\uplus L=\{1,\ldots, k\},$ $g'\in E$ such that $g_1\in E_1,$ $\supp(g')\cap \{j_1,\ldots, j_t\}=\emptyset,$ $g_2,\ldots, g_{2m}\in E_1$ and $\supp(g_i)\cap\supp(g_j)=\emptyset$ for all $i\neq j,$ we get $2^m \alpha_L g_1g_2\cdots g_{2m}e_{l_1}\cdots e_{l_s}=0 $ that is $\alpha_L=0,$ a contradiction. Therefore, $f=0$ and we are done in this case.

    Now suppose that $\deg f=2m+1$ and write
    $$
    f=  \sum_{L}\alpha_L[e_{l_1},x_{1}][x_2,x_3]\cdots [x_{2m},x_{2m+1}]e_{l_2}\cdots e_{l_s} \quad(\text{mod }\I^E(E^{(k)})),
    $$
    where $L=\{l_1,\ldots, l_s\}.$ With the same evaluation as before, one can get $f=0$ also in this case. 
    
    Thus these polynomials are linearly independent modulo $\I^E(E^{(k)})$ and since $P^E_n\cap I\subseteq P_n^E\cap\I^E(E^{(k)}),$ this proves that $I=\I^E(E^{(k)})$ and the polynomials in \eqref{polinomigradopari} and \eqref{polinomigradodispari} form a basis of $P_n^E$ modulo $P_n^E\cap\I^E(E^{(k)}).$

    If we set $\gamma_{2m}(E)$ and $\gamma_{2m+1}(E)$ to be the number of polynomials in \eqref{polinomigradopari} and \eqref{polinomigradodispari}, respectively, then it is clear that $\gamma_{2m}(E)=2^k,$ $\gamma_{2m+1}(E)= 2^k-1$ for all $m\geq 1$ and
    $$
    c_n^E(E^{(k)})=\sum_{t=0}^n\binom{n}{t}\gamma_t(E).
    $$
    Therefore by counting
    \begin{equation*}
    \begin{split}
        c_n^E(E^{(k)}) &= \sum_{h=0}^{\lfloor \frac{n}{2}\rfloor}\binom{n}{2h}2^k+\sum_{h=0}^{\lfloor \frac{n-1}{2}\rfloor}\binom{n}{2h+1}(2^k-1) = 2^k\left[\sum_{h=0}^{\lfloor \frac{n}{2}\rfloor}\binom{n}{2h}+\sum_{h=0}^{\lfloor \frac{n-1}{2}\rfloor}\binom{n}{2h+1}\right]- \sum_{h=0}^{\lfloor \frac{n-1}{2}\rfloor}\binom{n}{2h+1}%= 2^k2^n-2^{n-1}
        \\
        &= 2^{n-1}(2^{k+1}-1),
    \end{split}
    \end{equation*}
    as claimed.
\end{proof}

\end{document}
